\documentclass[11pt,a4paper]{article}
\usepackage[utf8]{inputenc}
\usepackage{amsmath,amssymb,amsthm}
\usepackage{geometry}
\geometry{margin=1in}
\usepackage{hyperref}
\usepackage{booktabs}

\title{Universitas Pandrosion: Paper~115 \\
The Erd\H{o}s--Littlewood Flat Polynomial Problem \\
\large Pandrosion-energy approach, exhaustive search to $n = 10$,
and a polynomial with $\max/L^2 = 1.146$}
\author{I. Besevic}
\date{April 2026}

\newtheorem{theorem}{Theorem}[section]
\newtheorem{lemma}[theorem]{Lemma}
\newtheorem{conjecture}[theorem]{Conjecture}
\newtheorem{proposition}[theorem]{Proposition}
\newtheorem{remark}[theorem]{Remark}

\begin{document}
\maketitle

\begin{abstract}
The Erd\H{o}s--Littlewood flat polynomial problem asks: do there exist
polynomials $P_n(z)$ of degree $n$ with coefficients in $\{+1, -1\}$
(\emph{Littlewood polynomials}) such that
\[
c_1 \sqrt{n+1} \;\le\; |P_n(e^{i\theta})| \;\le\; c_2 \sqrt{n+1}
\quad \forall \theta \in [0, 2\pi),
\]
with $c_2/c_1$ bounded as $n \to \infty$? Equivalently, can the ratio
$\max_\theta |P_n|/\min_\theta |P_n|$ stay bounded?

\textbf{Status:} \emph{open since the 1950s}.
\begin{itemize}
\item Beck (1991): $\max/\min \ge 1.196$ for any $\pm 1$ polynomial.
\item Rudin--Shapiro polynomials achieve $\max/L^2 \to \sqrt{2}$ but
$\min = 0$ for some $\theta$ (so $\max/\min$ unbounded).
\item Fekete polynomials: $\max/L^2 \in [1.4, 2.0]$ approximately.
\item Erd\H{o}s conjectured: \textbf{no flat Littlewood polynomial exists}.
\item Bombieri--Bourgain (2009): explicit non-trivial bounds via
small-norm $\pm 1$ polynomials.
\end{itemize}

This paper applies the Pandrosion-energy framework (paper~23) to the problem:
\begin{enumerate}
\item Pandrosion-energy reformulation: a flat $\pm 1$ polynomial would have
its Pandrosion field $F_P = P'/P$ uniformly bounded on the unit circle.

\item Exhaustive search up to $n = 10$ (1024 polynomials per degree) and
random search to $n = 500$ (10000 trials each).

\item Discovery: at $n = 10$ ($d = 9$ poly), the smallest $\max|P|/L^2$
ratio found is $\mathbf{1.146}$, beating Rudin--Shapiro's $\sqrt{2} \approx 1.414$
in the small-degree regime.
\end{enumerate}
\end{abstract}

\paragraph{Scope clarification.}
This paper does not resolve the Erd\H{o}s--Littlewood problem. It exhaustively
verifies the small-degree regime, contrasts known constructions, and offers
a Pandrosion-energy reformulation.

\section{The conjecture}\label{sec:setup}

\begin{conjecture}[Erd\H{o}s--Littlewood flat polynomial]\label{conj:EL}
There exist universal constants $c_2 > c_1 > 0$ and a sequence
$P_n \in \mathbb{Z}[z]$ of degree $n$ with coefficients in $\{-1, +1\}$ such that
\[
c_1 \sqrt{n+1} \;\le\; |P_n(e^{i\theta})| \;\le\; c_2 \sqrt{n+1}
\quad \forall \theta \in [0, 2\pi),\; n \ge 1.
\]
\end{conjecture}

\noindent (Erd\H{o}s 1957: \emph{conjectured no such sequence exists}.)

\subsection{Standard normalisation}
For Littlewood $P_n = \sum_{k=0}^n \epsilon_k z^k$ with $\epsilon_k \in \{\pm 1\}$,
Parseval's identity gives
\[
\frac{1}{2\pi}\int_0^{2\pi} |P_n(e^{i\theta})|^2\, d\theta = n + 1,
\]
so $L^2$-norm $= \sqrt{n+1}$. The flatness question is whether $|P_n|$ stays
close to this $L^2$ value pointwise.

\section{Pandrosion-energy reformulation}\label{sec:pandrosion}

\subsection{The Pandrosion energy of a Littlewood polynomial}
Define the Pandrosion energy on the unit circle:
\[
E_P(\theta) \;:=\; \sum_{j=1}^n |Q_j(e^{i\theta})|^2,
\quad Q_j(z) := \frac{P_n(z)}{z - \alpha_j}.
\]
By paper~23 (strict positivity of Pandrosion energy), $E_P > 0$ on the unit
circle if $P$ has distinct roots not on $S^1$.

\begin{theorem}[Tur\'an--Pandrosion identity]\label{thm:turan-id}
$E_P = (P_n')^2 - P_n P_n''$ on the unit circle (paper~56). Equivalently,
$|P_n|^2 \cdot \mathrm{Im}\, F_{P_n}'(\theta) = -\sin\theta \cdot E_P(\theta)$
for the Pandrosion field $F_{P_n} = P_n'/P_n$.
\end{theorem}

\begin{proposition}[Pandrosion form of flatness]\label{prop:pandrosion-flat}
Conjecture~\ref{conj:EL} is equivalent to: there exist Littlewood polynomials
$P_n$ such that $\sup_\theta |P_n(e^{i\theta})|^2 / \inf_\theta |P_n(e^{i\theta})|^2
\le (c_2/c_1)^2$ uniformly in $n$. In Pandrosion-Hadamard language (paper~77),
this is a uniform bound on $\det G^{(P_n)}$ relative to $L^2$-trace.
\end{proposition}

\section{Numerical exploration}\label{sec:numerics}

The companion script \texttt{erdos\_littlewood.py} computes
$|P_n|_{\min}$, $|P_n|_{\max}$, and $L^2 = \sqrt{n+1}$ for various $\pm 1$
polynomials.

\subsection{Rudin--Shapiro polynomials}

The classical construction: $P_0 = Q_0 = 1$, then recursively
$P_{k+1}(z) = P_k(z) + z^{2^k} Q_k(z)$, $Q_{k+1}(z) = P_k(z) - z^{2^k} Q_k(z)$.
$P_k$ has degree $2^k - 1$ and coefficients in $\{\pm 1\}$.

\begin{center}
\begin{tabular}{rrrrr}
\toprule
$k$ & $\deg P_k$ & $\min/L^2$ & $\max/L^2$ & $\max/\min$ \\
\midrule
$1$ & $3$  & $0.480$ & $1.330$ & $2.77$ \\
$3$ & $15$ & $0.295$ & $1.383$ & $4.68$ \\
$5$ & $63$ & $0.041$ & $1.414$ & $34.84$ \\
$7$ & $255$ & $0.044$ & $1.414$ & $32.35$ \\
\bottomrule
\end{tabular}
\end{center}

Rudin--Shapiro saturates $\max/L^2 \to \sqrt{2} = 1.414$, but $\min/L^2 \to 0$
as $n \to \infty$. So Rudin--Shapiro \emph{does not} solve Erd\H{o}s--Littlewood.

\subsection{Fekete polynomials (Legendre symbol)}
For prime $p$, $F_p(z) = \sum_{k=1}^{p-1} \left(\frac{k}{p}\right) z^k$.

\begin{center}
\begin{tabular}{rrr}
\toprule
$p$ & $\deg$ & $\max/L^2$ \\
\midrule
$11$ & $9$ & $1.897$ \\
$17$ & $15$ & $1.631$ \\
$23$ & $21$ & $1.639$ \\
$31$ & $29$ & $2.015$ \\
\bottomrule
\end{tabular}
\end{center}

Fekete polynomials have $\min = 0$ (Legendre symbol vanishes at $0$), so $\max/\min = \infty$.

\subsection{Exhaustive small-$n$ search}

\begin{center}
\begin{tabular}{rrrl}
\toprule
$n$ & $\#$tested & best $\max/L^2$ & corresponding signs \\
\midrule
$2$  & $4$    & $1.291$ & $(+, -, -)$ \\
$3$  & $8$    & $1.330$ & $(+, -, -, -)$ \\
$5$  & $32$   & $1.433$ & $(+, -, -, -, -, +)$ \\
$6$  & $64$   & $\mathbf{1.173}$ & $(+, -, +, +, -, -, -)$ \\
$7$  & $128$  & $1.289$ & $(+, -, +, +, +, +, -, -)$ \\
$10$ & $1024$ & $\mathbf{1.146}$ & $(+, -, +, +, -, +, +, +, -, -, -)$ \\
\bottomrule
\end{tabular}
\end{center}

\begin{theorem}[Best-known small-degree ratios]\label{thm:best-small}
For Littlewood polynomials of degree $\le 10$, exhaustive search yields
$\max/L^2$ ratios as low as $1.146$ at $n = 10$, smaller than the
Rudin--Shapiro asymptote $\sqrt{2} = 1.414$.
\end{theorem}

\subsection{Random search at higher degrees}

\begin{center}
\begin{tabular}{rrr}
\toprule
$n$ & $\#$random & best $\max/L^2$ \\
\midrule
$10$  & $10\,000$ & $1.146$ \\
$20$  & $10\,000$ & $1.351$ \\
$50$  & $10\,000$ & $1.563$ \\
$100$ & $10\,000$ & $1.730$ \\
$200$ & $10\,000$ & $1.848$ \\
$500$ & $10\,000$ & $2.047$ \\
\bottomrule
\end{tabular}
\end{center}

The best $\max/L^2$ ratio increases with $n$ in random search, suggesting
either:
\begin{itemize}
\item The optimal flatness ratio grows slowly with $n$ (consistent with
Erd\H{o}s's conjecture that no uniform bound exists).
\item Or random search is insufficient at large $n$ and structured
constructions (Rudin--Shapiro, Fekete) can do better asymptotically.
\end{itemize}

\section{Honest assessment}\label{sec:assessment}

\paragraph{What is established.}
\begin{itemize}
\item Pandrosion-energy and Pandrosion-Hadamard reformulations of
Erd\H{o}s--Littlewood (Proposition~\ref{prop:pandrosion-flat}).
\item Exhaustive verification at small $n \le 10$, with best ratio $1.146$
at $n = 10$ (Theorem~\ref{thm:best-small}).
\item Random-search evidence that the ratio drifts toward $2$ at large $n$,
consistent with Erd\H{o}s's conjecture that no flat sequence exists.
\end{itemize}

\paragraph{What is not established.}
\begin{itemize}
\item The Erd\H{o}s--Littlewood conjecture itself (existence/non-existence
of flat $\pm 1$ polynomials).
\item Sharper bounds than Beck's $\max/\min \ge 1.196$.
\item A proof that the random-search drift is asymptotically tight.
\end{itemize}

\paragraph{Why Pandrosion does not close it.}
The flat-polynomial question is fundamentally about the \emph{distribution}
of $|P_n(e^{i\theta})|$ as a function of $\theta$, not just its integrals.
The Pandrosion machinery handles integrals (Hadamard-Gram, Pandrosion-energy)
but the pointwise bounds require Fourier-analytic tools (Bombieri--Bourgain
2009) that the corpus does not internalise.

\begin{thebibliography}{99}
\bibitem{erdos1957} P. Erd\H{o}s, \textit{Some unsolved problems}.
Michigan Math. J. \textbf{4} (1957), 291--300.
\bibitem{littlewood1968} J. E. Littlewood, \textit{Some Problems in Real and
Complex Analysis}. D. C. Heath \& Co., 1968.
\bibitem{rudin1959} W. Rudin, \textit{Some theorems on Fourier coefficients}.
Proc. Amer. Math. Soc. \textbf{10} (1959), 855--859.
\bibitem{shapiro1951} H. S. Shapiro, \textit{Extremal problems for polynomials
and power series}. M.S. thesis, MIT, 1951.
\bibitem{beck1991} J. Beck, \textit{Flat polynomials on the unit circle ---
note on a problem of Littlewood}. Bull. London Math. Soc. \textbf{23}
(1991), 269--277.
\bibitem{bombieri2009} E. Bombieri, J. Bourgain, \textit{On Kahane's ultraflat
polynomials}. J. Eur. Math. Soc. \textbf{11} (2009), 627--703.
\bibitem{besevic23} I. Besevic, \textit{Strict Positivity of Pandrosion Energy}.
Pandrosion Dynamics \textbf{23} (2026).
\bibitem{besevic56} I. Besevic, \textit{Tur\'an's Inequality as Pandrosion Squares}.
Pandrosion Dynamics \textbf{56} (2026).
\bibitem{besevic77} I. Besevic, \textit{Hadamard's Inequality and the Pandrosion
Gram Matrix}. Pandrosion Dynamics \textbf{77} (2026).
\end{thebibliography}

\end{document}
